\documentclass{article}

% if you need to pass options to natbib, use, e.g.:
%     \PassOptionsToPackage{numbers, compress}{natbib}
% before loading neurips_2024


% ready for submission
\usepackage[final]{neurips_2024}


% to compile a preprint version, e.g., for submission to arXiv, add add the
% [preprint] option:
%     \usepackage[preprint]{neurips_2024}


% to compile a camera-ready version, add the [final] option, e.g.:
%     \usepackage[final]{neurips_2024}


% to avoid loading the natbib package, add option nonatbib:
%    \usepackage[nonatbib]{neurips_2024}


\usepackage[utf8]{inputenc} % allow utf-8 input
\usepackage[T1]{fontenc}    % use 8-bit T1 fonts
\usepackage{hyperref}       % hyperlinks
\usepackage{url}            % simple URL typesetting
\usepackage{booktabs}       % professional-quality tables
\usepackage{amsfonts}       % blackboard math symbols
\usepackage{nicefrac}       % compact symbols for 1/2, etc.
\usepackage{microtype}      % microtypography
\usepackage{xcolor}         % colors
\usepackage{amsmath}
\usepackage{graphicx}
\usepackage{subcaption}



\title{Nonparametric Density Estimation on High Dimensional Data}


% The \author macro works with any number of authors. There are two commands
% used to separate the names and addresses of multiple authors: \And and \AND.
%
% Using \And between authors leaves it to LaTeX to determine where to break the
% lines. Using \AND forces a line break at that point. So, if LaTeX puts 3 of 4
% authors names on the first line, and the last on the second line, try using
% \AND instead of \And before the third author name.

\author{%
Bryan Ng \\
University of Washington\\
Department of Statistics \\
Seattle, WA 98105 USA \\
\texttt{ngcw0303@uw.edu}\\
}

\begin{document}


\maketitle


\begin{abstract}
  Main reference paper: Wasserman, L., Liu, H., \& Lafferty, J. (2007). Nonparametric density estimation in high dimensions using the rodeo. Annals of Statistics, 35(4), 1765–1796.
  
  In this report, when considering the challenge of estimating the joint density function of a $d$-dimensional random vector 
  $X = (X_1, X_2, \ldots, X_d)$ in high-dimensional data space, we introduce a methodology called \emph{rodeo} (Regularization of the Derivative Expectation Operator) that simultaneously perform bandwidth selection and variable selection, and achieve faster rates of convergence compared to the kernel density estimator in high-dimensional setting.

\end{abstract}

\section{Background}

When addressing the challenge of estimating an unknown density, the kernel density estimator (KDE) is a widely-used method because it does not necessarily assuming that the underlying density from a specific parametric family. While KDE performs adequately when the dimension $d$ is relatively low, it encounters significant limitations as $d$ increased. The primary issues stem from the excessive computational demands of cross-validation for bandwidth selection in each dimension and the slow convergence rates of the estimators in high dimensions.


\subsection{Curse of Dimensionality}
Let $X_1, X_2, \ldots, X_n \in \mathbb{R}^d$ be a sample from a distribution $F$
with density $f$, and $x$ be a $d$-dimensional target point at which we want to estimate $f(x)$. The general multivariate kernel density estimator [1] is:

$$
\hat{f}_\mathbf{H}(x) = \frac{1}{n \det(\mathbf{H})} \sum_{i=1}^n \mathcal{K}\left(\mathbf{H}^{-1}(x - X_i)\right)
$$

where $\mathcal{K}$ is a symmetric kernel with $\int \mathcal{K}(u) du = 1$ and $\int u \mathcal{K}(u) du = 0_d$, $\mathbf{H} \in \mathbb{R}^{d \times d}$ is a symmetric positive definite bandwidth matrix. $\mathcal{K}$ is a product kernel and $\mathbf{H} = \text{diag}(h_1, \ldots, h_d)$ is diagonal, so that we can express the estimator as:

$$
\hat{f}_\mathbf{H}(x) = \frac{1}{n \det(\mathbf{H})} \sum_{i=1}^n \mathcal{K}\left(\mathbf{H}^{-1}(x - X_i)\right)
$$

$$
= \frac{1}{n} \sum_{i=1}^n \prod_{j=1}^d \frac{1}{h_j} \mathcal{K}\left(\frac{x_j - X_{ij}}{h_j}\right).
$$

The multivariate kernel density estimator has a slower convergence rate [2]:

$$
\text{MISE}(\hat{f_{\mathbf{H}}}) = O(h^4) + O\left(\frac{1}{nh^d}\right) \rightarrow \text{MISE}_{\text{opt}}(\hat{f_{\mathbf{H}}}) = O\left(n^{-\frac{4}{4+d}}\right)
$$
which highlights the "curse of dimensionality" when $d$ is large.

\subsection{Computational Cost of Bandwidth Selection}

The most extended criteria in bandwidth estimation is the asymptotical mean integrated square error (AMISE) [2]. In the one-dimensional case, optimizing this criteria with respect to the bandwidth is not problematic as it involves a global search over one variable which is computationally feasible.

$$\text{AMISE(h)} = \frac{\mu_K^2}{4} \cdot h^4 \cdot \int \left| p''(x) \right|^2 dx + \frac{\sigma_K^2}{n h}$$ 
The optimal smoothing bandwidth is often determined by minimizing this expression, leading to:
$$h_{\text{opt}} = \left( \frac{1}{n} \cdot \frac{4}{\int \left| p''(x) \right|^2 dx} \cdot \frac{\sigma_K^2}{\mu_K^2} \right)^{1/5}
$$

In the d-dimensional case, however, the AMISE becomes more complex [3]:
$$
\text{AMISE}(\mathbf{H}) = \frac{1}{4} \mu_2(K)^2 \left( \text{vech}^\top \mathbf{H} \right) \mathbf{\Psi}_4 \left( \text{vech} \, \mathbf{H} \right) 
+ n^{-1} \lvert \mathbf{H} \rvert^{-1/2} R(K).
$$

$$
\mathbf{H}_{\text{opt}} = \arg\min_\mathbf{H} \text{AMISE}(\mathbf{H})
$$
Solving the above optimization problem would be extremely difficult. However, we can derive $\mathbf{H}_{\text{opt}}$ through unbiased cross-validation as an alternative way. Namely, the optimal bandwidth matrix $\mathbf{H}_{\text{opt}}$ is obtained by minimizing the UCV criterion [3]:

$$
\text{UCV}(\mathbf{H}) = \left( \int \hat{f}_\mathbf{H}^2(x) \, dx - \frac{2}{n} \sum_{i=1}^n \hat{f}_{\mathbf{H},-i}(X_i) \right)
$$

$$
\mathbf{H}_{\text{opt}} = \arg\min_\mathbf{H} \text{UCV}(\mathbf{H}) = \arg\min_\mathbf{H} \left( \int \hat{f}_\mathbf{H}^2(x) \, dx - \frac{2}{n} \sum_{i=1}^n \hat{f}_{\mathbf{H},-i}(X_i) \right)
$$

But it is also impractical due to high computational costs of efficiently computing the integrals and sums.

Consequently, kernel density estimators are infrequently applied to multivariate data because of the challenges in determining suitable kernel bandwidths to avoid over-fitting as well as the relatively slow rate of convergence.

\section{Rodeo Density Estimation}

\subsection{Sparsity Assumption}

To effectively navigate density estimation, it is essential to define a suitable sparsity condition. Our primary assumption is:
$$
f(x_1, \dots, x_d) = g(x_R) b(x_1, \dots, x_d)
$$
where $g$ is an unknown function, $x_R = (x_j : j \in R)$ is a subset of the variables, and $b$ is a baseline density that is either fully known or constrained by a finite number of parameters. If $R$ contains only a few coordinates, it follows that the nonparametric component $g$ depends on a limited number of variables. Consequently, the rodeo method can achieve a faster convergence rate of 

$$
O\left( n^{-\frac{4}{4+R}} \right)
$$

where $R \ll d$.


\subsection{The Local Rodeo}

Considering the kernel density estimator $\hat{f}_H(x)$, the density rodeo is based on the concept of searching for $H_{\text{opt}}$. We initiate with a bandwidth matrix $H = \text{diag}(h_0, \ldots, h_0)$, where $h_0$ is large. Then, we compute the derivatives $Z_j$ for $j = 1, \ldots, d$ of the kernel density estimate concerning each bandwidth $h_j$, reducing $h_j$ if $Z_j$ is significant. The test statistic is defined as:

$$
Z_j = \frac{\partial \hat{f}_{\mathbf{H}(x)}}{\partial h_j} = \frac{1}{n} \sum_{i=1}^{n} \frac{\partial}{\partial h_j} \prod_{k=1}^{d} \frac{1}{h_k} K\left( \frac{x_k - X_{ik}}{h_k} \right)
$$

Thus, $|Z_j|$ is large if altering $h_j$ significantly impacts the estimator. To conduct the test, we compare $Z_j$ with its variance:

$$
\sigma^2_j = \text{Var}(Z_j) = \text{Var}\left(\frac{1}{n} \sum_{i=1}^{n} Z_{ji}\right) = \frac{1}{n} \text{Var}(Z_{j1}).
$$

For estimation, we use $s^2_j = \frac{v^2_j}{n}$ where $v^2_j$ denotes the sample variance of the $Z_j$ values. For a general kernel, we have:

$$
Z_j = -\frac{1}{n} \sum_{i=1}^{n} \frac{1}{h_j^2} (x_j - X_{ij})^2 \tilde{K}\left( \frac{x_j - X_{ij}}{h_j} \right) \prod_{k=1}^{d} K\left( \frac{x_k - X_{ik}}{h_k} \right),
$$

where $\tilde{K}(x) = \frac{d}{dx} \log K(x)$. In the case where $K$ is the Gaussian kernel, this simplifies to:

$$
Z_j = C \cdot \frac{1}{n} \sum_{i=1}^{n} (x_j - X_{ij})^2 - h_j^2 \cdot K\left( \frac{x_k - X_{ik}}{h_k} \right).
$$

Here, the constant $C = \frac{1}{h_j^3} \prod_{k=1}^{d} \frac{1}{h_k}$ can be disregarded to prevent overflow in calculations as $h_k \to 0$ for large $d$. The test statistic $Z_j$ indicates dimensions where changes in the bandwidth $h_j$ significantly influence the density estimate, necessitating greater precision in the kernel, thus requiring smaller $h_j$. Consequently, the local rodeo effectively accomplishes bandwidth selection.


\subsection{The Global Rodeo}

The local rodeo can be adapted for global bandwidth selection, where each dimension utilizes a fixed bandwidth. This strategy averages test statistics across multiple evaluation points $(x_1, \ldots, x_k)$, sampled from the empirical distribution.

Averaging $Z_j$ yields a statistic with an asymptotic mean for relevant variables of:

$$
\frac{1}{h_j} \sum_{k=1}^{k} f_{jj}(x_i)
$$

Due to potential sign changes in $f_{jj}(x)$, cancellations may result, leading to a small statistic. To mitigate this, the statistic is squared. Let $x_1, \ldots, x_m$ be the evaluation points, and $Z_j(x_i)$ denotes the derivative for the $i$-th evaluation point with respect to $h_j$:

$$
Z_j(x_i) = \frac{1}{n} \sum_{k=1}^{n} Z_{jk}(x_i), \quad i = 1, \ldots, m, \quad j = 1, \ldots, d
$$

Define $\gamma_{jk} = (Z_{j1}(x_k), Z_{j2}(x_k), \ldots, Z_{jm}(x_k))^T$. Assuming $\text{Var}(\gamma_{jk}) = \Sigma_j$, let $Z_j$ denote $ (Z_{j1}, Z_{j2}, \ldots, Z_{jm})^T$. By the multivariate central limit theorem, we find:

$$
\text{Var}(Z_j) = \frac{\Sigma_j}{n} \equiv C_j
$$

Thus, we define the test statistic as:

$$
T_j = \frac{1}{m} \sum_{k=1}^{m} Z^2_j(x_k), \quad j = 1, \ldots, d
$$

and the variance estimate:

$$
s_j = \sqrt{\text{Var}(T_j)} = \frac{1}{m} \sqrt{\text{Var}(Z^T_j Z_j)} = \frac{1}{m} \left(2 \text{tr}(C_j^2) + 4 \mu^T_j C_j \mu_j\right)
$$

where $\hat{\mu} = \frac{1}{m} \sum_{i=1}^{m} Z_j(x_i)$. For irrelevant dimensions $j \in R_c$, we have $E[Z_j(x_i)] = O(h_j)$, thus $E[T_j] \approx \text{Var}(Z_j(x_i))$. We estimate variance using $s^2_j$. The threshold is set as:

$$
\lambda_j = s^2_j + 2 s_j \sqrt{2 \log\left(\frac{n}{c_n}\right)}
$$



\subsection{The algorithm of LDER(Local Density Estimation Rodeo)}

\begin{enumerate}
    \item Select parameter $0 < \beta < 1$ and initial bandwidth $h_0 = \frac{c_0}{\log \log n}$ for some constant $c_0$. Also, let $c_n = O(\log n)$.
    \item Initialize $h_j$, and activate all dimensions:
    \begin{enumerate}
        \item $h_j = h_0, \, j = 1, 2, \ldots, d.$
        \item $\mathcal{A} = \{1, 2, \ldots, d\}.$
    \end{enumerate}
    \item While $\mathcal{A}$ is nonempty, do for each $j \in \mathcal{A}$:
    \begin{enumerate}
        \item Estimate the derivative and variance: $Z_j$ and $s_j^2$.
        \item Compute the threshold $\lambda_j = s_j \sqrt{2 \log (nc_n)}.$
        \item If $\lvert Z_j \rvert > \lambda_j$, then set $h_j \leftarrow \beta h_j$; otherwise remove $j$ from $\mathcal{A}$.
    \end{enumerate}
    \item Output bandwidths $\mathbf{H}^*$ and estimator $\hat{f}_{\mathbf{H}^*}(x)$ .
\end{enumerate}

\section{Example}

We demonstrate example of density estimation of two-dimensional example by rodeo and kernel density estimation respectively [4]. The two dimensions are independently generated from $X_1$, $X_2$ with the sample size $n = 500$, where
$$
X_1 \sim \frac{2}{3}\text{Beta}(1, 2) + \frac{1}{3}\text{Beta}(10, 10)
$$

$$
X_2 \sim \text{Uniform}(0, 1)
$$

\begin{figure}[h!]
    \centering
    % First image
    \begin{subfigure}[b]{0.45\textwidth}
        \centering
        \includegraphics[width=\textwidth]{rodeo.png}
        \caption{the rodeo estimation}
        \label{fig:sub1}
    \end{subfigure}
    \hfill
    % Second image
    \begin{subfigure}[b]{0.45\textwidth}
        \centering
        \includegraphics[width=\textwidth]{KDE.png}
        \caption{the KDE estimation}
        \label{fig:sub2}
    \end{subfigure}
    \caption{Density estimation on two-dimensional dataset}
    \label{fig:parallel_images}
\end{figure}

We observe that the results obtained using the rodeo estimation outperform those of kernel density estimation.

\section{Conclusion}

The rodeo methodology provides a robust strategy for tackling the challenges of nonparametric kernel density estimation in high-dimensional spaces. By effectively selecting bandwidths and identifying relevant variables, it considerably reduces the curse of dimensionality along with the computational costs associated with bandwidth optimization for each individual dimension.

\section{References}

{
\small

[1] Scott, D. W. (2015). Multivariate density estimation: Theory, practice, and visualization (2nd ed.). Wiley. https://doi.org/10.1002/9781118575574

[2] Wasserman, L. (2006). All of nonparametric statistics. Springer. https://doi.org/10.1007/0-387-30623-4

[3] Duong, T., \& Hazelton, M. L. (2005). Cross-validation bandwidth matrices for multivariate kernel density estimation. Scandinavian Journal of Statistics, 32(3), 485–506. https://doi.org/10.1111/j.1467-9469.2005.00445.x

[4] Wasserman, L., Liu, H., \& Lafferty, J. (2007). Nonparametric density estimation in high dimensions using the rodeo. Annals of Statistics, 35(4), 1765–1796.



}

\end{document}